\chapter{Bine Butterfly Algorithms}\label{ch:butterfly}

% Channel colours used by the figures of this chapter
\colorlet{pch0}{ch0!70}
\colorlet{pch1}{ch1!70}
\colorlet{pch2}{ch2!70}
\colorlet{pch3}{chunkB!90}

This chapter implements the Bine butterfly of \cref{sec:bg-bine-butterfly} in \nccl{} for the allgather, the reduce-scatter and the allreduce.
\Cref{sec:bfly-impl} describes the schedule of the distance-doubling butterfly, \cref{sec:bfly-collectives} how it is turned into kernels, \cref{sec:bfly-dhlv} the distance-halving schedule used by the double send strategy, and \cref{sec:bfly-testing} the invariants that the schedule satisfies.

%==============================================================================
\section{Schedule construction}\label{sec:bfly-impl}

As for the tree, the schedule is built for $p=2^s$ ranks and is defined only for powers of two.
Unlike the tree, a butterfly has no root, so there is no rotation: one set of tables serves every collective on the communicator.

\subsection{Output and notation}

The schedule consists of a table and two vectors,
\begin{align}
  \Pi(r,i) &\in \{0,\dots,p-1\}, \\
  \nu(r) &\in \{0,\dots,p-1\}, \qquad \nu^{-1}(x)\in\{0,\dots,p-1\},
\end{align}
for $r\in[0,p-1]$ and $i\in[0,s-1]$.
$\Pi(r,i)=q$ means that rank $r$ exchanges data with rank $q$ at step $i$, and $\nu$ is a rank-dependent label (\cref{sec:bfly-def}) with inverse $\nu^{-1}$.
All the ranks communicate at all the steps, so the table has no empty entries, unlike the tables of the tree (\cref{sec:tree-impl}).
The arrays are filled by
\begin{verbatim}
void ncclGetBineButterflyDdbl(int nRanks, int steps,
                              int *partners, int *index, int *order);
\end{verbatim}
which writes \texttt{partners} (the table $\Pi$), \texttt{index} (the vector $\nu$) and \texttt{order} (the vector $\nu^{-1}$).
Step $i$ is the $i$-th step of the reduce-scatter, where the distance between partners grows; the allgather is its reverse.

\subsection{Definition}\label{sec:bfly-def}

Let
\begin{equation}
  \rho_i=\sum_{k=0}^{i}(-2)^k=\frac{1-(-2)^{i+1}}{3}.
  \label{eq:bfly-rho}
\end{equation}
The partner of rank $r$ at step $i$ is
\begin{equation}
  \Pi(r,i)=
  \begin{cases}
    (r+\rho_i)\bmod p & \text{if $r$ is even},\\
    (r-\rho_i)\bmod p & \text{if $r$ is odd}.
  \end{cases}
  \label{eq:bfly-partner}
\end{equation}
This is the partner of \cref{eq:bg-bine-bfly-partner} with the steps in the opposite order, since $\rho_i=\tau_{s-1-i}$ (\cref{eq:bg-bine-tau}), that is, $\Pi(r,i)=\pi_{s-1-i}(r)$.
The distance $\lvert\rho_i\rvert\approx 2^{i+1}/3$ grows with $i$.

The label $\nu$ is the position of a rank in the distance-doubling \bine{} tree rooted at rank $0$~\cite{desensi2025bine}.
With
\begin{equation}
  h(r)=
  \begin{cases}
    0 & \text{if } r=0,\\
    \rtn(p-r) & \text{if $r$ is even and } r\neq0,\\
    \rtn(r) & \text{if $r$ is odd,}
  \end{cases}
  \qquad
  \nu(r)=h(r)\oplus\bigl(h(r)\gg1\bigr),
  \label{eq:bfly-nu}
\end{equation}
where $\gg$ is the logical right shift, and $\nu^{-1}$ is its inverse (it exists by invariant~\ref{inv:bfly-perm}).
The labels define where the block of a rank is stored in the kernels: the \emph{position} of the block of rank $r$ is
\begin{equation}
  \mathrm{pos}(r)=\mathrm{rev}(\nu(r)),
\end{equation}
where $\mathrm{rev}$ reverses the $s$ bits.
We will see that the blocks exchanged at a step are contiguous in the position space and, in general, not in the rank space.

\subsection{Example: $p=8$}

For $p=8$ ($s=3$), $\rho_0=1$, $\rho_1=-1$ and $\rho_2=3$.
\Cref{tab:bfly-nu-p8} gives the labels and the positions, and \cref{fig:bfly-butterfly} the partners.
For example, rank $3$ is odd and communicates with $3-1=2$, $3+1=4$ and $3-3=0$.

\begin{table}[t]
\centering
\small
\caption{Values $h(r)$ and $\nu(r)$ (binary) and position $\mathrm{pos}(r)=\mathrm{rev}(\nu(r))$ (decimal) for $p=8$. The inverse is $\nu^{-1}=[0,1,7,2,3,6,4,5]$.}
\label{tab:bfly-nu-p8}
\begin{tabular}{@{}l*{8}{c}@{}}
\toprule
$r$ & 0 & 1 & 2 & 3 & 4 & 5 & 6 & 7 \\
\midrule
$h(r)$ & 000 & 001 & 010 & 111 & 100 & 101 & 110 & 011 \\
$\nu(r)$ & 000 & 001 & 011 & 100 & 110 & 111 & 101 & 010 \\
$\mathrm{pos}(r)$ & 0 & 4 & 6 & 1 & 3 & 7 & 5 & 2 \\
\bottomrule
\end{tabular}
\end{table}

\begin{figure}[t]
\centering
\begin{tikzpicture}[x=2.3cm, y=-0.62cm, font=\footnotesize,
    rk/.style={circle, draw, fill=white, minimum size=0.6cm, inner sep=0pt}]
  % partner edges: column i -> column i+1
  \foreach \r/\q in {0/1,1/0,2/3,3/2,4/5,5/4,6/7,7/6}
    \draw[ch0, thick] (0,\r) -- (1,\q);
  \foreach \r/\q in {0/7,1/2,2/1,3/4,4/3,5/6,6/5,7/0}
    \draw[ch1, thick] (1,\r) -- (2,\q);
  \foreach \r/\q in {0/3,1/6,2/5,3/0,4/7,5/2,6/1,7/4}
    \draw[ch2, thick] (2,\r) -- (3,\q);
  % nodes
  \foreach \c in {0,...,3}
    \foreach \r in {0,...,7}
      \node[rk] at (\c,\r) {\r};
  % labels
  \node[ch0] at (0.5,-1.1) {$i=0$, $\lvert\rho_0\rvert=1$};
  \node[ch1] at (1.5,-1.1) {$i=1$, $\lvert\rho_1\rvert=1$};
  \node[ch2] at (2.5,-1.1) {$i=2$, $\lvert\rho_2\rvert=3$};
\end{tikzpicture}
\caption{Distance-doubling \bine{} butterfly for $p=8$. An edge from rank $r$ in column $i$ to rank $q$ in column $i+1$ means $\Pi(r,i)=q$. The reduce-scatter reads the figure from left to right, the allgather from right to left.}
\label{fig:bfly-butterfly}
\end{figure}

\subsection{Computing the tables}

\Cref{alg:bfly-schedule} computes the three arrays.
The labels are computed first and $\nu^{-1}$ is obtained by inverting them; then the partners are computed one step at a time, accumulating $\rho_i$ in a 64-bit integer and reducing it modulo $p$ with a non-negative result.
The cost is $O(p\log p)$ time and memory.

\begin{algorithm}[t]
\caption{Schedule of the distance-doubling Bine butterfly ($p=2^s$).}
\label{alg:bfly-schedule}
\begin{algorithmic}[1]
\Require $p=2^s$
\Ensure table $\Pi$ of size $p\times s$, vectors $\nu$ and $\nu^{-1}$ of size $p$
\For{$r\gets 0$ \textbf{to} $p-1$}
  \State compute $h(r)$ \Comment{\cref{eq:bfly-nu}}
  \State $\nu(r)\gets h(r)\oplus(h(r)\gg1)$;\quad $\nu^{-1}(\nu(r))\gets r$
\EndFor
\State $\rho\gets0$;\quad $w\gets1$
\For{$i\gets 0$ \textbf{to} $s-1$}
  \State $\rho\gets\rho+w$;\quad $w\gets-2w$ \Comment{$\rho=\rho_i$}
  \State $m\gets\rho\bmod p$ \Comment{non-negative}
  \For{$r\gets 0$ \textbf{to} $p-1$}
    \State $\Pi(r,i)\gets (r+m)\bmod p$ \textbf{if} $r$ is even, \textbf{else} $(r-m)\bmod p$
  \EndFor
\EndFor
\end{algorithmic}
\end{algorithm}

%==============================================================================
\section{Allgather, reduce-scatter and allreduce}\label{sec:bfly-collectives}

The table says who communicates with whom; which data moves is decided by the collective.
\begin{itemize}
  \item \emph{Allgather.} Each rank contributes a block of \texttt{count} elements and obtains the $p$ blocks, the block of rank $b$ at offset $b\cdot\texttt{count}$ of \texttt{recvbuff}.
  \item \emph{Reduce-scatter.} Each rank holds $p$ blocks, and rank $b$ obtains the reduction of the blocks $b$ of all the ranks.
  \item \emph{Allreduce.} Each rank obtains the elementwise reduction of the vectors of all the ranks; it is a reduce-scatter followed by an allgather (\cref{sec:bfly-allreduce}).
\end{itemize}

\subsection{Integration in \nccl{}}

The butterfly is integrated as the tree (\cref{sec:tree-integration}), with the same algorithm \texttt{NCCL\_ALGO\_BINE} and pattern \texttt{ncclPatternBine}, and three pieces.
\begin{itemize}
  \item \emph{Pattern.} \texttt{ncclPatternBine}, used by the host to compute the number of steps and the chunking.
  \item \emph{Tables.} The tables of \cref{sec:bfly-impl,sec:bfly-dhlv} are stored in \texttt{channel.bine} next to those of the tree: \texttt{partners} ($\Pi$), \texttt{dhlvPartners} ($\Pi^{\mathrm{dh}}$), \texttt{index} ($\nu$) and \texttt{order} ($\nu^{-1}$).
  The fields \texttt{nSteps} and \texttt{nDoublingSteps} are both set to $s$, and the kernels use the latter.
  \item \emph{\texttt{RunWorkColl} specializations.} One for each collective and protocol, as for the tree, which call the kernel of the strategy of \cref{sec:bfly-strategies} (send, permute, block-by-block or double send).
\end{itemize}

\subsection{Blocks, positions and strategies}\label{sec:bfly-strategies}

\paragraph{Allgather as the reverse of the reduce-scatter.}
Step $k\in[0,s-1]$ of the allgather uses the partner $\Pi(r,s-1-k)$: the distance decreases and the data doubles at each step.
At step $k$ a rank holds $g=2^k$ blocks, sends them and receives the $g$ blocks of its partner.
The reduce-scatter is the time reverse: step $i$ uses $\Pi(r,i)$, and the data halves.

\paragraph{Contiguity in the position space.}
Let $x=\mathrm{pos}(r)$ and $y=\mathrm{pos}(q)$ for $q=\Pi(r,s-1-k)$.
Bit $b$ of $\nu$ is bit $s-1-b$ of $\mathrm{pos}$, so invariant~\ref{inv:bfly-split} with $i=s-1-k$ says that $x$ and $y$ agree above bit $k$ and differ at bit $k$.
Hence the $g$ blocks held at step $k$ are those whose positions lie in the aligned range of size $g$ containing $x$, and the blocks of the partner lie in the sibling range (\cref{fig:bfly-held}).
In the position space the pattern is a recursive-doubling allgather; in the rank space, the blocks of an aligned range are the ranks $\nu^{-1}(\mathrm{rev}(y))$, which are in general not contiguous.

\begin{figure}[t]
\centering
\begin{tikzpicture}[x=1.05cm, y=0.95cm, font=\footnotesize,
    cell/.style={draw, minimum width=1cm, minimum height=0.8cm, inner sep=0pt, align=center}]
  \foreach \row/\n in {0/1,1/2,2/4,3/8}
    \foreach \y/\b in {0/0,1/3,2/7,3/4,4/1,5/6,6/2,7/5} {
      \ifnum\y<\n \def\fc{ch0!35}\else\def\fc{white}\fi
      \node[cell, fill=\fc] at (\y,-\row) {$\y$\\[-1pt]{\tiny $r{=}\b$}};
    }
  \node[anchor=east] at (-0.7,0)  {start};
  \node[anchor=east] at (-0.7,-1) {after $k{=}0$};
  \node[anchor=east] at (-0.7,-2) {after $k{=}1$};
  \node[anchor=east] at (-0.7,-3) {after $k{=}2$};
  \node[anchor=west] at (7.7,-1) {$q=3$ (pos.\ 1)};
  \node[anchor=west] at (7.7,-2) {$q=7$ (pos.\ 2)};
  \node[anchor=west] at (7.7,-3) {$q=1$ (pos.\ 4)};
\end{tikzpicture}
\caption{Blocks held by rank $0$ (position $0$) during the allgather with $p=8$. Cell $y$ is a position; $r$ is the rank whose block it contains, $\nu^{-1}(\mathrm{rev}(y))$. The held positions (shaded) always form an aligned range, whereas the held ranks $\{0\}$, $\{0,3\}$, $\{0,3,7,4\}$ are not contiguous. The right column is the partner at each step.}
\label{fig:bfly-held}
\end{figure}

\paragraph{Strategies.}
The strategies differ in where they put the blocks, not in the pattern.
Each is defined by the \emph{slot} of position $y$ (the block-sized region of the working buffer used for it), a \emph{pre-step} that puts the data in the right slots, and a \emph{post-step} that restores the natural order (\cref{tab:bfly-strategies}).
\begin{itemize}
  \item \emph{Send} exchanges the blocks before the butterfly: $r$ sends its block to $\nu^{-1}(\mathrm{rev}(r))$ and receives that of rank $\mathrm{pos}(r)$, stored at slot $\mathrm{pos}(r)$.
  Ranks with $\mathrm{pos}(r)=r$ need no exchange.
  \item \emph{Permute} stores the block at slot $\mathrm{pos}(r)$ and permutes the blocks in place at the end.
  \item \emph{Block-by-block} keeps the natural order, so the slot of position $y$ is $\nu^{-1}(\mathrm{rev}(y))$ and every block is sent separately.
  \item \emph{Double send} keeps the natural order with the distance-halving schedule, whose windows are contiguous modulo $p$ (\cref{sec:bfly-dhlv}).
\end{itemize}
In the reduce-scatter the roles of exchange and permutation are mirrored: the exchange of \emph{send} becomes a post-step, and the permutation of \emph{permute} a pre-step.
When a reduce-scatter is followed by an allgather, as in the allreduce, the exchange of \emph{send} cancels out~\cite{desensi2025bine}.

\paragraph{Fused copies.}
The copy of the own block is not a separate pass in the allgather: it is fused into the first send (\texttt{copySend}), which reads \texttt{sendbuff}, writes the slot and sends.
It is executed separately, followed by a barrier, only when it cannot be fused: in place, where a direct write of the partner could overwrite the source.
In the reduce-scatter the last receive writes the reduced block directly in \texttt{recvbuff} (except in \emph{send}, which is followed by the exchange).

\begin{table}[t]
\centering
\footnotesize
\setlength{\tabcolsep}{4pt}
\caption{The strategies as choices of slot, pre-step and post-step. Send, permute and block-by-block use the distance-doubling table (reversed in the allgather); double send uses the distance-halving one. Copies are fused as described in the text.}
\label{tab:bfly-strategies}
\begin{tabularx}{\linewidth}{@{}l l X X l@{}}
\toprule
Strategy & Slot of $y$ & Allgather & Reduce-scatter & Unit \\
\midrule
Send & $y$ & pre: exchange with $\nu^{-1}(\mathrm{rev}(r))$ & post: same exchange & run \\
Permute & $y$ & pre: copy to slot $\mathrm{pos}(r)$; post: permutation & pre: permutation; post: copy slot $\mathrm{pos}(r)$ & run \\
Block-by-block & $\nu^{-1}(\mathrm{rev}(y))$ & pre: copy to slot $r$ & post: copy slot $r$ & block \\
Double send & circular window & pre: copy to slot $r$ & post: copy slot $r$ & $\le2$ runs \\
\bottomrule
\end{tabularx}
\end{table}

\subsection{Work partition}\label{sec:bfly-pos}

By default a channel handles a slice of \emph{every} block (\emph{slice mode}).
A range of $g$ positions is then $g$ pieces separated by a stride of \texttt{count}, contiguous only if \texttt{channelCount}$=$\texttt{count}.
For small messages this multiplies the number of calls by the number of channels.

Send, double send and block-by-block therefore have a second partition, \emph{pos mode}, in which a channel handles \emph{whole} blocks of a range of positions (\cref{fig:bfly-partition}): with $n_c$ channels, channel $j$ owns the positions $[j\lceil p/n_c\rceil,(j+1)\lceil p/n_c\rceil)$, clipped to $p$.
The part of a window owned by a channel is then one contiguous run for the strategies whose slots are the positions.
Pos mode is selected when
\begin{equation}
  n_c>1,\qquad p\ge n_c,\qquad \texttt{count}\cdot\texttt{sizeof}(T)<n_c\cdot S,
  \label{eq:bfly-pos}
\end{equation}
where $S$ is a threshold (\texttt{BINE\_POS\_SLICE\_BYTES}).
In other words, a block is partitioned by position only when its slice per channel, $\texttt{count}\cdot\texttt{sizeof}(T)/n_c$, would be smaller than $S$, and the block size at which the partition changes is $n_c S$.
The condition depends only on values equal on all ranks and channels.

The value of $S$ was determined experimentally: we swept it and measured, for each value, the latency of the strategies in both partitions.
The two partitions cross at about $S=16\,\text{KiB}$ per channel, which is the default value.
$S$ is a compile-time parameter: \texttt{BINE\_POS\_SLICE\_BYTES} is defined with this default only if it is not already defined, so it can be changed at build time (e.g., with \texttt{-DBINE\_POS\_SLICE\_BYTES=}$\langle\text{bytes}\rangle$) without touching the source, and the crossover can differ on other systems.
Below it pos mode wins, since a contiguous run needs fewer calls than the strided pieces of slice mode (\cref{sec:bfly-res-modes}); above it the idle channels of pos mode cost more than the calls it saves, so slice mode wins.

\begin{figure}[t]
\centering
\begin{tikzpicture}[x=0.95cm, y=0.9cm, font=\footnotesize]
  % legend
  \foreach \c in {0,...,3} {
    \fill[pch\c] (2.2*\c,1.9) rectangle ++(0.3,0.3);
    \node[anchor=west] at (2.2*\c+0.3,2.05) {channel \c};
  }
  % (a) slice mode
  \node[anchor=south west] at (0,0.75) {(a) slice mode};
  \foreach \y in {0,...,7} {
    \foreach \c in {0,...,3}
      \fill[pch\c] (\y+0.25*\c,0) rectangle ++(0.25,0.6);
    \draw (\y,0) rectangle ++(1,0.6);
    \node[below] at (\y+0.5,0) {\y};
  }
  \draw[decorate, decoration={brace, amplitude=4pt}] (8,-0.65) -- (4,-0.65)
    node[midway, below=4pt, align=center] {window of 4 positions:\\every channel moves 4 strided pieces};
  % (b) pos mode
  \begin{scope}[yshift=-4.3cm]
    \node[anchor=south west] at (0,0.75) {(b) pos mode};
    \foreach \y/\c in {0/0,1/0,2/1,3/1,4/2,5/2,6/3,7/3} {
      \fill[pch\c] (\y,0) rectangle ++(1,0.6);
      \draw (\y,0) rectangle ++(1,0.6);
      \node[below] at (\y+0.5,0) {\y};
    }
    \draw[decorate, decoration={brace, amplitude=4pt}] (8,-0.65) -- (4,-0.65)
      node[midway, below=4pt, align=center] {window of 4 positions:\\channels 2 and 3 move one contiguous run each};
  \end{scope}
\end{tikzpicture}
\caption{Work partition for $p=8$ blocks and four channels. In slice mode every block is cut among the channels; in pos mode each channel owns two whole blocks.}
\label{fig:bfly-partition}
\end{figure}

Pos mode needs no synchronization among channels because it is \emph{channel-closed}: a position is written, and re-sent, only by the channel that owns it, and the range that a rank sends is the range that its partner receives, so both ends clip it with the same ownership and issue the same transfers.
A channel with an empty intersection does not take part in that direction of the step.
Permute cannot use pos mode: its final permutation moves the block of rank $b$ from position $\mathrm{pos}(b)$ to $b$, across the ranges of different channels, and \nccl{} has no barrier among channels.

\subsection{Kernels}\label{sec:bfly-kernels}

\paragraph{Structure.}
\Cref{alg:bfly-allgather,alg:bfly-reducescatter} show the common structure of the kernels of send, permute and block-by-block.
They differ only in \textsc{PreStep}, \textsc{Slot}, \textsc{PostStep} (\cref{tab:bfly-strategies}) and in the partition.
$O$ is the set of positions owned by the channel ($[0,p)$ in slice mode) and $n_s$ the number of threads of the send group.

\begin{algorithm}[t]
\caption{Bine allgather, kernel of one channel.}
\label{alg:bfly-allgather}
\begin{algorithmic}[1]
\State \Call{PreStep}{}\Comment{send: exchange; otherwise copy, fused into step $0$}
\State $x\gets\mathrm{pos}(r)$
\For{$k \gets 0$ \textbf{to} $s-1$}
  \State $q\gets\Pi(r,s-1-k)$;\quad $g\gets2^k$;\quad $b\gets\lfloor x/2g\rfloor\cdot2g$
  \If{bit $k$ of $x$ is $0$} $(a_s,a_r)\gets(b,\,b+g)$ \textbf{else} $(a_s,a_r)\gets(b+g,\,b)$ \EndIf
  \If{$\mathit{tid}<n_s$}\Comment{send group}
    \For{each transfer of the slots \Call{Slot}{$y$}, $y\in[a_s,a_s+g)\cap O$}
      \State \texttt{copySend} to $q$ if $k=0$ and fused, else \texttt{sendFromOutput}
    \EndFor
  \Else\Comment{receive group}
    \For{each transfer of the slots \Call{Slot}{$y$}, $y\in[a_r,a_r+g)\cap O$}
      \State \texttt{recv} from $q$
    \EndFor
  \EndIf
  \State \textbf{barrier}
\EndFor
\State \Call{PostStep}{}\Comment{permute: in-place permutation}
\end{algorithmic}
\end{algorithm}

\begin{algorithm}[t]
\caption{Bine reduce-scatter, kernel of one channel.}
\label{alg:bfly-reducescatter}
\begin{algorithmic}[1]
\State \Call{PreStep}{}\Comment{permute: in-place permutation of the input, then barrier}
\State $x\gets\mathrm{pos}(r)$
\For{$i \gets 0$ \textbf{to} $s-1$}
  \State $q\gets\Pi(r,i)$;\quad $g\gets2^{s-1-i}$;\quad $b\gets\lfloor x/2g\rfloor\cdot2g$
  \If{bit $s-1-i$ of $x$ is $0$} $(a_s,a_r)\gets(b+g,\,b)$ \textbf{else} $(a_s,a_r)\gets(b,\,b+g)$ \EndIf
  \If{$\mathit{tid}<n_s$}\Comment{send group}
    \For{each transfer of the slots \Call{Slot}{$y$}, $y\in[a_s,a_s+g)\cap O$}
      \State \texttt{sendFromOutput} to $q$
    \EndFor
  \Else\Comment{receive group}
    \For{each transfer of the slots \Call{Slot}{$y$}, $y\in[a_r,a_r+g)\cap O$}
      \State \texttt{recvReduceCopy} from $q$ \Comment{in place; the last step applies the final operation}
    \EndFor
  \EndIf
  \State \textbf{barrier}
\EndFor
\State \Call{PostStep}{}\Comment{send: exchange; otherwise copy to the output if not fused}
\end{algorithmic}
\end{algorithm}

\paragraph{Separate send and receive groups.}
In the butterfly a rank sends and receives with the same peer at every step, and the two directions are independent.
The first $n_s=\lfloor n/2/32\rfloor\cdot32$ threads of the channel therefore send, with fan $\langle0,1\rangle$, and the others receive, with fan $\langle1,0\rangle$, each with its own \texttt{Primitives} (\cref{tab:bfly-calls}).
A send that waits for a FIFO slot thus never blocks the receive group.
Three points follow.
\begin{itemize}
  \item A group builds its \texttt{Primitives} only if it has traffic, and both ends of a connection must agree on that: the constructor of the direct variants exchanges pointers with the peer, and a lone constructor would wait forever.
  \item The receive group has its own buffers, so that at the last step it can write the result in \texttt{recvbuff} while the send group still reads the working buffer.
  \item A barrier separates the steps, since the next step sends, and in the reduce-scatter reduces, what the previous one wrote.
\end{itemize}
The in-place permutation follows the cycles of the permutation, so no scratch buffer is needed.
In the reduce-scatter it completes before any \texttt{Primitives} is built, since a direct peer may read the buffer.

\begin{table}[t]
\centering
\footnotesize
\setlength{\tabcolsep}{4pt}
\caption{\nccl{} operations of the allgather (AG) and reduce-scatter (RS). The direct operation is used when the zero-copy path is available.}
\label{tab:bfly-calls}
\begin{tabularx}{\linewidth}{@{}l X c l l@{}}
\toprule
 & Role & Fan & Direct & Otherwise \\
\midrule
AG & Send the window & $\langle0,1\rangle$ & \texttt{directSendFromOutput} & \texttt{sendFromOutput} \\
 & First send, own block fused & $\langle0,1\rangle$ & \texttt{directCopySend} & \texttt{copySend} \\
 & Receive the window & $\langle1,0\rangle$ & \texttt{directRecv} & \texttt{recv} \\
RS & Send the window & $\langle0,1\rangle$ & \texttt{directSendFromOutput} & \texttt{sendFromOutput} \\
 & Receive and reduce the window & $\langle1,0\rangle$ & \texttt{directRecvReduceCopy} & \texttt{recvReduceCopy} \\
Both & Exchange of \emph{send} & $\langle0,1\rangle$, $\langle1,0\rangle$ & \texttt{directSend}, \texttt{directRecv} & \texttt{send}, \texttt{recv} \\
\bottomrule
\end{tabularx}
\end{table}

\paragraph{Transfers.}\label{sec:bfly-fast}
No transfer is larger than one chunk (\texttt{chunkCount} elements), so that each call corresponds to one step of the slot FIFO (\cref{sec:nccl-transport}).
How a window is cut depends on its layout in memory.
\begin{itemize}
  \item \emph{Contiguous run.} If \texttt{channelCount}$=$\texttt{count} (one channel, or pos mode) and the slots are adjacent, the part of the window owned by the channel is one region of $n_B\cdot\texttt{count}$ elements, sent as $\lceil n_B\cdot\texttt{count}/\texttt{chunkCount}\rceil$ chunks, possibly across block boundaries.
  This is the case of send and permute; in double send a window is at most two runs (\cref{sec:bfly-dhlv}).
  \item \emph{Strided transfers.} Otherwise the kernel issues one transfer per chunk and per block, $n_B\lceil\texttt{channelCount}/\texttt{chunkCount}\rceil$ in total.
  This is always the case of block-by-block, and of every strategy in slice mode when the slice does not cover whole blocks.
\end{itemize}

\subsection{Contract with the proxy}\label{sec:bfly-proxy}

The proxy must post as many FIFO steps as the kernel consumes, per channel, step and direction.
Since every call consumes the FIFO steps of one chunk whatever the number of elements it moves, empty calls included (\cref{sec:nccl-fifo-cost}), the proxy only needs the number of calls, provided that no call moves more than one chunk (\cref{sec:bfly-fast}).
With $n_B$ the number of blocks of a window owned by the channel, $c$ the chunk size and $c_c=\texttt{channelCount}$, a send or a receive of a window costs:
\begin{itemize}
  \item $n_B\lceil c_c/c\rceil$ calls in slice mode with $c_c<\texttt{count}$, and for block-by-block in every mode;
  \item $\lceil n_B\,\texttt{count}/c\rceil$ calls per contiguous run otherwise (summed over the at most two runs of a circular window).
\end{itemize}
The two directions are counted separately, since in pos mode a channel may send and not receive at a step.
The exchange of \emph{send} costs $\lceil c_c/c\rceil$ calls in each direction on every channel in slice mode, and $\lceil\texttt{count}/c\rceil$ on the single channel that owns the position in pos mode.
The proxy evaluates \cref{eq:bfly-pos} with the same formula as the kernel, so that both derive the same partition and the same number of calls.

\subsection{Allreduce}\label{sec:bfly-allreduce}

The allreduce is the reduce-scatter of \cref{alg:bfly-reducescatter} followed by the allgather of \cref{alg:bfly-allgather} on the same butterfly.
The result is the whole vector, so there is no block to restore: after the reduce-scatter a rank holds the reduced block of position $\mathrm{pos}(r)$, and the allgather returns every block to its position.
The exchange of \emph{send} cancels out, so the allreduce has no pre-step and no post-step and uses \emph{send} without the exchange.

\paragraph{Layout.}
The slice of the channel is processed in $L=\lceil\texttt{channelCount}/(pc)\rceil$ loops of $p$ chunks of $c$ elements; block $b$ of loop $l$ is the chunk at offset $\texttt{grid}+l\,pc+b\,c$.
Chunks beyond the end of the slice are shorter or empty, and the calls are still issued so that both ends of an exchange match.
The kernel first copies \texttt{sendbuff} into \texttt{recvbuff} when they differ, then works in place in \texttt{recvbuff}.

\paragraph{Chunking.}
The chunk is limited to one FIFO step, so every call sends and receives a single chunk and the kernel needs no contiguous runs.
The price is $g\cdot L$ sends and as many receives at a step with a window of $g$ blocks.
The proxy must derive the same chunk size.

\paragraph{Kernel.}
\Cref{alg:bfly-allreduce} shows the kernel; a rank skips a step when its partner is undefined or is itself, and the proxy applies the same condition.
\todo{The allreduce kernel is not refactored yet: it still uses one \texttt{Primitives} with fan $\langle1,1\rangle$ per step instead of the send and receive groups of \cref{sec:bfly-kernels}.}

\begin{algorithm}[t]
\caption{Bine allreduce, kernel of one channel.}
\label{alg:bfly-allreduce}
\begin{algorithmic}[1]
\State copy \texttt{sendbuff} to \texttt{recvbuff} if they differ;\ \textbf{barrier}
\State $x\gets\mathrm{pos}(r)$
\For{$i \gets 0$ \textbf{to} $s-1$}\Comment{reduce-scatter}
  \State $q\gets\Pi(r,i)$;\quad $g\gets2^{s-1-i}$;\quad $b\gets\lfloor x/2g\rfloor\cdot2g$
  \If{bit $s-1-i$ of $x$ is $0$} $(a_s,a_r)\gets(b+g,\,b)$ \textbf{else} $(a_s,a_r)\gets(b,\,b+g)$ \EndIf
  \For{$l\gets0$ \textbf{to} $L-1$ and $j\gets0$ \textbf{to} $g-1$}
    \State \texttt{sendFromOutput} to $q$ the chunk $a_s+j$ of loop $l$
    \State \texttt{recvReduceCopy} from $q$ the chunk $a_r+j$ of loop $l$
  \EndFor
  \State \textbf{barrier}
\EndFor
\For{$k \gets 0$ \textbf{to} $s-1$}\Comment{allgather}
  \State $q\gets\Pi(r,s-1-k)$;\quad $g\gets2^k$;\quad $b\gets\lfloor x/2g\rfloor\cdot2g$
  \If{bit $k$ of $x$ is $0$} $(a_s,a_r)\gets(b,\,b+g)$ \textbf{else} $(a_s,a_r)\gets(b+g,\,b)$ \EndIf
  \For{$l\gets0$ \textbf{to} $L-1$ and $j\gets0$ \textbf{to} $g-1$}
    \State \texttt{sendFromOutput} to $q$ the chunk $a_s+j$ of loop $l$
    \State \texttt{recv} from $q$ the chunk $a_r+j$ of loop $l$
  \EndFor
  \State \textbf{barrier}
\EndFor
\end{algorithmic}
\end{algorithm}

%==============================================================================
\section{Distance-halving schedule (double send)}\label{sec:bfly-dhlv}

The double send strategy runs the butterfly with the distance-halving schedule of \cref{eq:bg-bine-bfly-partner} in the order in which it is defined.
The table is filled by
\begin{verbatim}
void ncclGetBineButterflyDhlv(int nRanks, int steps, int *partners);
\end{verbatim}
which takes no root and produces no labels: the blocks stay in the natural rank order, so $\nu$ and $\nu^{-1}$ are not needed.

\subsection{Definition}

Let $\tau_i=\sum_{j=0}^{s-i-1}(-2)^j=\frac{1-(-2)^{s-i}}{3}$, the offset of \cref{eq:bg-bine-bfly-partner}.
The partner of rank $r$ at step $i$ is
\begin{equation}
  \Pi^{\mathrm{dh}}(r,i)=
  \begin{cases}
    (r+\tau_i)\bmod p & \text{if $r$ is even},\\
    (r-\tau_i)\bmod p & \text{if $r$ is odd}.
  \end{cases}
  \label{eq:bfly-dhlv-partner}
\end{equation}
The distance $\lvert\tau_i\rvert\approx2^{s-i}/3$ halves at each step.
Since $\rho_i=\tau_{s-1-i}$, the table is that of \cref{sec:bfly-impl} with the steps reversed, $\Pi^{\mathrm{dh}}(r,i)=\Pi(r,s-1-i)$: the two generators differ only in the order of the steps.
For $p=8$, $\tau_0=3$, $\tau_1=-1$ and $\tau_2=1$, so the table is \cref{fig:bfly-butterfly} read from right to left.

\subsection{Circular window}

In this schedule the blocks exchanged are contiguous in the rank space if the buffer is considered circular (\cref{sec:bg-bine-butterfly}).
In the reduce-scatter, before step $i$ a rank holds the contributions to a window of $2^{s-i}$ consecutive blocks modulo $p$, sends one half to its partner and keeps the other.
With $g=s-1-i$, each half has $2^g$ blocks, and the two partners hold the same window and keep complementary halves.
The rank $r$ keeps the upper half if and only if $r$ and $g$ have different parity, so it sends the lower half in that case and the upper half otherwise.
The window start moves up by $2^g$ at the steps where the upper half is kept, and the last window is $\{r\}$, so the first window starts at
\begin{equation}
  a_0(r)=\Bigl(r-\sum_{\substack{0\le g<s\\ g\not\equiv r\ (\mathrm{mod}\ 2)}}2^g\Bigr)\bmod p .
  \label{eq:bfly-dhlv-start}
\end{equation}
\Cref{fig:bfly-circular} shows the windows of rank $0$ for $p=8$: $a_0=6$, and the first window wraps around the end of the buffer.
At step $0$ rank $0$ sends to rank $3$ the blocks $[2,5]$, which are those that rank $3$ keeps.

\begin{figure}[t]
\centering
\begin{tikzpicture}[x=1.05cm, y=0.95cm, font=\footnotesize,
    cell/.style={draw, minimum width=1cm, minimum height=0.8cm, inner sep=0pt}]
  \foreach \row/\sent/\kept in {0/{2,3,4,5}/{6,7,0,1}, 1/{6,7}/{0,1}, 2/{1}/{0}} {
    \foreach \y in {0,...,7} \node[cell] at (\y,-\row) {\y};
    \foreach \y in \sent \node[cell, fill=crosslink!35] at (\y,-\row) {\y};
    \foreach \y in \kept \node[cell, fill=ch2!35] at (\y,-\row) {\y};
  }
  \node[anchor=east] at (-0.7,0)  {$i=0$};
  \node[anchor=east] at (-0.7,-1) {$i=1$};
  \node[anchor=east] at (-0.7,-2) {$i=2$};
  \node[anchor=west] at (7.7,0)  {to rank 3};
  \node[anchor=west] at (7.7,-1) {to rank 7};
  \node[anchor=west] at (7.7,-2) {to rank 1};
  % legend
  \node[cell, fill=crosslink!35, minimum width=0.5cm, minimum height=0.4cm] at (1,-3.1) {};
  \node[anchor=west] at (1.4,-3.1) {sent};
  \node[cell, fill=ch2!35, minimum width=0.5cm, minimum height=0.4cm] at (3.5,-3.1) {};
  \node[anchor=west] at (3.9,-3.1) {kept};
  \node[cell, minimum width=0.5cm, minimum height=0.4cm] at (6,-3.1) {};
  \node[anchor=west] at (6.4,-3.1) {outside the window};
\end{tikzpicture}
\caption{Reduce-scatter windows of rank $0$ with the distance-halving schedule, $p=8$. The window of step $0$ starts at block $6$ and wraps around: the kept half is $\{6,7,0,1\}$. At each step the kept half is the window of the next step, and the last is $\{0\}$.}
\label{fig:bfly-circular}
\end{figure}

The allgather is the time reverse: it starts from $\{r\}$ and doubles the window at each step, reading the table from the last step to the first.
It therefore moves the largest data over the longest distance, the opposite of the other strategies; this is the price of keeping the natural order without any pre-step or post-step.

\subsection{Kernels}

The kernels have the structure of \cref{alg:bfly-allgather,alg:bfly-reducescatter} with $\Pi^{\mathrm{dh}}$ and the circular window in place of the position space.
Position and block coincide, so both partitions of \cref{sec:bfly-pos} apply.
The only copy is that of the own block, fused as in \cref{sec:bfly-strategies}.
A window of $2^g$ blocks starting at $a$ is at most two runs, $[a,\min(a+2^g,p))$ and, if it wraps, $[0,a+2^g-p)$, each intersected with the positions owned by the channel and cut into chunks when \texttt{channelCount}$=$\texttt{count}, and block by block otherwise.
Sender and receiver visit the runs in the same order, since the window received is the one the partner sends.
The window start is updated at every step, also when the rank does not take part in it.

%==============================================================================
\section{Testing}\label{sec:bfly-testing}

Both schedules are deterministic functions of $p=2^s$, and every property the collectives rely on is a statement about finite tables.
We therefore test them exhaustively against a set of invariants, for every $p=2^s$ with $1\le s\le 12$ (up to $2^{12}=4{,}096$ ranks), every rank and every step.
There is no root to enumerate, so the cost is $O(p\log p)$ lookups per value of $p$, instead of the $O(p^2 s)$ of the tree.
The tests run on the host and cover the two generators separately: \texttt{ncclGetBineButterflyDdbl}, which produces $\Pi$, $\nu$ and $\nu^{-1}$ (\cref{sec:bfly-impl}), and \texttt{ncclGetBineButterflyDhlv}, which produces $\Pi^{\mathrm{dh}}$ (\cref{sec:bfly-dhlv}).

\subsection{Invariants}

\paragraph{Distance-doubling table.}
For every $p=2^s$, rank $r$ and step $i\in[0,s-1]$:
\begingroup
\renewcommand{\theenumi}{J\arabic{enumi}}
\renewcommand{\labelenumi}{\textbf{(\theenumi)}}
\begin{enumerate}
  \item \label{inv:bfly-sym} \emph{Symmetry.} $\Pi(\Pi(r,i),i)=r$.
  \item \label{inv:bfly-self} \emph{No self-communication.} $\Pi(r,i)\neq r$.
  \item \label{inv:bfly-full} \emph{Full participation.} $\Pi(r,i)\in[0,p-1]$ is defined for every $r$ and $i$.
  \item \label{inv:bfly-perm} \emph{Labels are a permutation.} $\nu$ is a permutation of $[0,p-1]$ and $\nu^{-1}(\nu(r))=r$.
  \item \label{inv:bfly-split} \emph{Partners split the blocks.} With $q=\Pi(r,i)$,
  \begin{equation}
    \bigl(\nu(r)\oplus\nu(q)\bigr)\bmod 2^{i+1}=2^i,
    \label{eq:bfly-split}
  \end{equation}
  i.e., the labels of $r$ and $q$ agree on the bits below $i$ and differ at bit $i$.
  \item \label{inv:bfly-closed} \emph{Closed form.} $\Pi(r,i)$ is given by \cref{eq:bfly-partner}.
\end{enumerate}
\endgroup

\paragraph{Distance-halving table.}
Let $\tau_i=(1-(-2)^{s-i})/3$ as in \cref{sec:bfly-dhlv}, and let $d(r,q)=\min\bigl((q-r)\bmod p,\,(r-q)\bmod p\bigr)$ be the modular distance.
For every $p=2^s$, rank $r$ and step $i\in[0,s-1]$:
\begingroup
\renewcommand{\theenumi}{K\arabic{enumi}}
\renewcommand{\labelenumi}{\textbf{(\theenumi)}}
\begin{enumerate}
  \item \label{inv:dh-sym} \emph{Symmetry.} $\Pi^{\mathrm{dh}}(\Pi^{\mathrm{dh}}(r,i),i)=r$.
  \item \label{inv:dh-self} \emph{No self-communication.} $\Pi^{\mathrm{dh}}(r,i)\neq r$.
  \item \label{inv:dh-full} \emph{Full participation.} $\Pi^{\mathrm{dh}}(r,i)\in[0,p-1]$ for every $r$ and $i$ (the generator initialises the table to $-1$, so an entry that is never written is detected).
  \item \label{inv:dh-parity} \emph{Opposite parity.} $\Pi^{\mathrm{dh}}(r,i)$ and $r$ have different parity.
  \item \label{inv:dh-closed} \emph{Closed form.} $\Pi^{\mathrm{dh}}(r,i)$ is given by \cref{eq:bfly-dhlv-partner}, with $\tau_i$ computed from the closed form above.
  \item \label{inv:dh-rev} \emph{Reversal.} $\Pi^{\mathrm{dh}}(r,i)=\Pi(r,s-1-i)$.
  \item \label{inv:dh-dist} \emph{Modular distance.} $d\bigl(r,\Pi^{\mathrm{dh}}(r,i)\bigr)=\lvert\tau_i\rvert$, and it does not increase with $i$.
  \item \label{inv:dh-window} \emph{Circular windows.} Let $a_0(r)$ be given by \cref{eq:bfly-dhlv-start}, $g=s-1-i$, and let the window of $r$ before step $i$ be the $2^{g+1}$ blocks starting at $a_i(r)$ (modulo $p$).
  The rank keeps the upper half if $r$ and $g$ have different parity, and the lower half otherwise, so $a_{i+1}(r)=(a_i(r)+2^g)\bmod p$ in the first case and $a_{i+1}(r)=a_i(r)$ in the second.
  Then, for $q=\Pi^{\mathrm{dh}}(r,i)$:
  \begin{enumerate}
    \item[(a)] $a_i(q)=a_i(r)$: the partners hold the same window;
    \item[(b)] the half sent by $r$ starts at $a_{i+1}(q)$: it is the half kept by the partner;
    \item[(c)] after the last step, $a_s(r)=r$: the final window is $\{r\}$.
  \end{enumerate}
\end{enumerate}
\endgroup

\subsection{Why they hold and what they guarantee}

\paragraph{Exchange structure (\ref{inv:bfly-sym})--(\ref{inv:bfly-full}), (\ref{inv:dh-sym})--(\ref{inv:dh-parity}).}
Both tables use offsets that are always odd: $\rho_i=\rho_{i-1}+(-2)^i$ with $\rho_0=1$, and $\tau_i$ is the sum of $1$ and of even terms, so they stay odd modulo the even number $p$.
An even rank adds the offset and an odd rank subtracts it, hence the partner of an even rank is odd, the partner of an odd rank is even (\ref{inv:dh-parity}), and no rank can be its own partner (\ref{inv:bfly-self}, \ref{inv:dh-self}).
If $r$ is even and $q=r+m$, then $q$ is odd and its partner is $q-m=r$; the case of $r$ odd is symmetric.
Hence the pairing is an involution (\ref{inv:bfly-sym}, \ref{inv:dh-sym}).
Both formulas are defined for every rank, so no entry is empty (\ref{inv:bfly-full}, \ref{inv:dh-full}).
At each step the ranks form $p/2$ pairs, sends and receives always match, and the kernels need no branch for ranks that skip a step, unlike the tree (\cref{sec:tree-testing}).

\paragraph{Labels (\ref{inv:bfly-perm}).}
The conversion $\rtn$ maps a rank to an $s$-bit negabinary string by value modulo $p$.
It is a bijection: two different strings, whose lowest differing bit is $j<s$, have values that differ by $2^j$ modulo $2^{j+1}$, hence modulo $p$.
The least significant bit of $\rtn(x)$ is the parity of $x$.
Odd ranks take $\rtn(r)$ and are mapped to $p/2$ distinct strings ending in $1$; even ranks take $\rtn(p-r)$, with $r=0$ taking the zero string, and are mapped to the $p/2$ distinct strings ending in $0$.
So $h$ is a bijection onto the $s$-bit strings, and so is $x\mapsto x\oplus(x\gg1)$, which is the Gray code.
Therefore no block is lost or stored twice, and $\nu^{-1}$ is well defined.

\paragraph{Correct exchange (\ref{inv:bfly-split}).}
This is the invariant that ties the partners to the blocks.
It relates two independent computations, the partners (closed form with the sums $\rho_i$) and the labels (Gray code of the negabinary representation), and it is not implied by the other invariants.
It says that two partners hold the same blocks and each keeps the half the other sends, so after $s$ steps every contribution is combined exactly once.
In the position space it says that partners agree above bit $k=s-1-i$ and differ at bit $k$, which makes the exchanged blocks contiguous (\cref{sec:bfly-strategies}) and pos mode channel-closed (\cref{sec:bfly-pos}).
The stronger statement $\nu(r)\oplus\nu(q)=2^i$ holds only for the tree rooted at rank $0$, which is why the test checks the bits up to $i$ only.

\paragraph{Identifying the tables (\ref{inv:bfly-closed}, \ref{inv:dh-closed}, \ref{inv:dh-rev}).}
Symmetry, absence of self-communication and full participation also hold for the standard recursive-doubling butterfly, and, for the distance-halving table, also for the table with the two signs exchanged between even and odd ranks.
The closed forms are what fixes the partners (and, for $\Pi^{\mathrm{dh}}$, the sign), and the reversal (\ref{inv:dh-rev}) ties the two generators together: since $\rho_i=\tau_{s-1-i}$, the distance-halving table must be the distance-doubling one read from the last step to the first.
The two generators are separate code, so (\ref{inv:dh-rev}) also transfers to $\Pi^{\mathrm{dh}}$ what was checked on $\Pi$.

\paragraph{Locality (\ref{inv:dh-dist}).}
Since $\lvert\tau_i\rvert=(2^{s-i}\pm1)/3\le p/2$, the modular distance between partners is exactly $\lvert\tau_i\rvert$, about $2/3$ of the $2^{s-i-1}$ of a standard distance-halving butterfly (Eq.~\eqref{eq:bg-bine-bfly-partner} and~\cite{desensi2025bine}).
The sequence $\lvert\tau_i\rvert$ is $3,1,1$ for $p=8$ and does not increase in general.
This is the property that gives Bine its advantage, and it does not depend on the sign, so it does not replace (\ref{inv:dh-closed}).

\paragraph{Windows (\ref{inv:dh-window}).}
This is the semantic check of the distance-halving table as a schedule of the reduce-scatter.
Before step $i$ a rank holds the contributions to a window of $2^{g+1}$ consecutive blocks modulo $p$, and sends one half to its partner while keeping the other (\cref{sec:bfly-dhlv}).
Two partners have opposite parity (\ref{inv:dh-parity}), so one keeps the upper half and the other the lower half, and if they hold the same window (a) the half that each sends is the half that the other keeps (b).
Condition (a) is not implied by the previous ones: at step $i+1$ the new partner of $r$ is given by the table, and it holds the same window as $r$ only if the table and the start $a_0(r)$ of \cref{eq:bfly-dhlv-start} are consistent.
By induction the window of a rank is halved at each step, from $p$ blocks (the whole buffer, starting at $a_0(r)$) down to one block, which by (c) is block $r$.
Hence every contribution to a block is combined exactly once and the block ends at its owner, and the allgather, its time reverse, starts from $\{r\}$ and ends with the whole buffer.
In the test the three conditions are checked for every rank and step by simulating the windows with the rule above.

\subsection{Scope and limits}

The tests are exhaustive over all ranks and steps for every power of two $p\le 2^{12}$, with a cost of $O(p\log p)$ lookups per value of $p$ and per test.
Beyond this range, correctness relies on the arguments of this section, not on testing.
They check the schedule only: the kernels, the work partition and the splitting of the circular windows into runs (\cref{sec:bfly-kernels}) are not covered here.
The oracles are of different strength.
(\ref{inv:bfly-closed}) recomputes $\rho_i$ with the same recurrence as the generator, so it only checks that the table realizes its defining formula, while (\ref{inv:dh-closed}) uses the closed form of $\tau_i$ and therefore an independent computation.
(\ref{inv:bfly-split}) and (\ref{inv:dh-window}) are the checks of the construction as a whole.
(\ref{inv:dh-window}) uses the keep rule and the start $a_0(r)$ of \cref{sec:bfly-dhlv} as oracle, so it validates the table against these rules, and the rules themselves rest on the argument above.

%==============================================================================
\section{Experimental results}\label{sec:bfly-results}

% Two stacked panels (allgather above, reduce-scatter below) at a fixed pixel scale,
% so that all plots of this section are printed with the same text size.
% \bflypair{allgather file}{reduce-scatter file}{width as a fraction of \linewidth}
\newcommand{\bflypair}[3]{%
  \centering
  \captionsetup[subfigure]{skip=2pt, font=small}%
  \begin{subfigure}[t]{\linewidth}
    \centering
    \includegraphics[width=#3\linewidth]{#1}
    \caption{Allgather.}
  \end{subfigure}\par\vspace{6pt}
  \begin{subfigure}[t]{\linewidth}
    \centering
    \includegraphics[width=#3\linewidth]{#2}
    \caption{Reduce-scatter.}
  \end{subfigure}}

This section compares the \bine{} butterfly allgather and reduce-scatter of \cref{sec:bfly-collectives} with the native \nccl{} ring (\cref{ch:nccl}), under the \texttt{SIMPLE} and \texttt{LL} protocols.
We first compare \bine{}, with its best strategy at each point, with the ring (\cref{sec:bfly-res-ring}), and then the four strategies with one another (\cref{sec:bfly-res-modes}).

\paragraph{Setup.}
Experiments run on the Leonardo supercomputer~\cite{cineca2024leonardo}, on nodes with four NVIDIA A100 GPUs connected by NVLink, with $p\in\{4,8,16,32,64\}$ ranks (one to 16 nodes, one rank per GPU) and block sizes from 1\,kB to 64\,MB (the block contributed by each rank in the allgather and obtained by each rank in the reduce-scatter, as in \cref{ch:tree}).
Each run has $10$ warmup and $40$ measured iterations, and each point is the median latency of the measured iterations, as in \cref{ch:treeallreduce}.
The variability across iterations is large for small messages (the median relative standard deviation is about $30\,\%$ at 1\,kB, $10\,\%$ at 256\,kB and $2\,\%$ at 64\,MB), so differences of a few percent at small and medium sizes are not significant.
Since a node holds four GPUs, $p=4$ is entirely intra-node.
\texttt{LL128} is not measured, for the reasons of \cref{sec:tree-results}.

%------------------------------------------------------------------------------
\subsection{\bine{} versus ring}\label{sec:bfly-res-ring}

\Cref{fig:bfly-scaling} shows the scaling with $p$ of the ring and of the best \bine{} strategy, \cref{fig:bfly-heatmap-simple,fig:bfly-heatmap-best} the fastest algorithm at each point (with \texttt{SIMPLE} only, and with the best protocol of each algorithm), and \cref{tab:bfly-best} the ratio between the ring and \bine{}, each with its best protocol.
The two collectives behave in the same way.

\begin{figure}[!tbp]
  \bflypair{AllGather_SIMPLE+LL_vs_ranks_by_size_bv-BUTTERFLY_bm-best.png}{ReduceScatter_SIMPLE+LL_vs_ranks_by_size_bv-BUTTERFLY_bm-best.png}{0.92}
  \caption{Median latency versus number of ranks, \bine{} butterfly with its best strategy (blue) against the \nccl{} ring (red), for six block sizes. Solid lines with filled markers: \texttt{SIMPLE}; dashed lines with hollow markers: \texttt{LL}. Both axes are logarithmic.}
  \label{fig:bfly-scaling}
\end{figure}

In this implementation \bine{} wins only for small messages on many ranks.
With \texttt{SIMPLE} on both sides it is faster than the ring in $8$ of the $30$ points of the allgather and $7$ of the reduce-scatter (1 and 16\,kB from $16$ and $8$ ranks respectively, plus 256\,kB on $64$ ranks in the allgather; \cref{fig:bfly-heatmap-simple}), but once the ring may use \texttt{LL} it is the fastest in only $3$ of the $60$ points (\cref{fig:bfly-heatmap-best}): at 1\,kB and $p=64$ it is $1.17\times$ (allgather) and $1.9\times$ (reduce-scatter) faster than the best ring, and $2.9$--$3.1\times$ faster than the ring with \texttt{SIMPLE}.

\begin{figure}[!tbp]
  \bflypair{AllGather_SIMPLE_best_algo_heatmap_bv-BUTTERFLY_bm-best.png}{ReduceScatter_SIMPLE_best_algo_heatmap_bv-BUTTERFLY_bm-best.png}{0.73}
  \caption{Fastest algorithm for each block size and number of ranks with the \texttt{SIMPLE} protocol (median latency). Each cell lists the other combinations with the percentage by which they are slower. Bf-S: \bine{} butterfly with its best strategy, R-S: \nccl{} ring.}
  \label{fig:bfly-heatmap-simple}
\end{figure}

\begin{table}[!tbp]
\centering
\small
\setlength{\tabcolsep}{4pt}
\renewcommand{\arraystretch}{1.1}
\sisetup{table-format=1.2, detect-weight=true, mode=text}
\caption{Median latency of the ring divided by that of the best \bine{} strategy, each with its best protocol, for the allgather and the reduce-scatter, with $p$ ranks. A value below $1$ means that the ring is faster; values above $1$ are in bold.}
\label{tab:bfly-best}
\begin{tabular}{@{}l *{5}{S} @{\hspace{2em}} *{5}{S}@{}}
\toprule
 & \multicolumn{5}{c}{Allgather} & \multicolumn{5}{c}{Reduce-scatter} \\
\cmidrule(lr){2-6}\cmidrule(l){7-11}
Block size & {$p{=}4$} & {8} & {16} & {32} & {64} & {$p{=}4$} & {8} & {16} & {32} & {64} \\
\midrule
1\,kB & 0.62 & 0.67 & 0.74 & 0.90 & {\bfseries 1.17} & 0.64 & 0.66 & 0.73 & 0.98 & {\bfseries 1.90} \\
16\,kB & 0.70 & 0.52 & 0.49 & 0.50 & 0.64 & 0.63 & 0.59 & 0.53 & 0.48 & 0.59 \\
256\,kB & 0.60 & 0.89 & 0.93 & 0.94 & {\bfseries 1.01} & 0.63 & 0.89 & 0.90 & 0.94 & 0.98 \\
1\,MB & 0.68 & 0.75 & 0.72 & 0.69 & 0.68 & 0.65 & 0.65 & 0.64 & 0.65 & 0.64 \\
4\,MB & 0.57 & 0.64 & 0.67 & 0.64 & 0.63 & 0.53 & 0.65 & 0.63 & 0.60 & 0.65 \\
64\,MB & 0.44 & 0.75 & 0.65 & 0.61 & 0.56 & 0.43 & 0.72 & 0.63 & 0.65 & 0.63 \\
\bottomrule
\end{tabular}
\end{table}

\begin{figure}[!tbp]
  \bflypair{AllGather_SIMPLE+LL_best_algo_heatmap_bv-BUTTERFLY_bm-best.png}{ReduceScatter_SIMPLE+LL_best_algo_heatmap_bv-BUTTERFLY_bm-best.png}{0.95}
  \caption{As \cref{fig:bfly-heatmap-simple}, with both protocols. The suffixes S and LL are the protocols.}
  \label{fig:bfly-heatmap-best}
\end{figure}

The reason is the one of \cref{sec:bg-ring-vs-tree} (solid lines of \cref{fig:bfly-scaling}): at 1\,kB the ring time doubles at every doubling of $p$ ($36$ to $556$\,\textmu s in the allgather from $4$ to $64$ ranks, with \texttt{SIMPLE}), while the \bine{} time grows by a constant $31$--$42$\,\textmu s ($48$ to $191$\,\textmu s).
The other point is the allgather at 256\,kB on $64$ ranks, where the two are on par (within $11\,\%$ for $p\ge8$ at this size).
At 16\,kB the ring with \texttt{LL} is $1.4$--$2.0\times$ faster at every scale.
From 1\,MB on the ring is faster by $1.3$--$1.8\times$ with \texttt{SIMPLE}, and the gap does not depend on $p$.
Both algorithms have the same bandwidth term, $\tfrac{p-1}{p}n\beta$ (\cref{tab:bg-cost}), so the difference is not in the volume of data but in how efficiently it is moved.
The main cause is the cost of a call (\cref{sec:nccl-fifo-cost}): every call consumes the FIFO steps of a full chunk, whatever the number of elements it moves.
The ring moves full chunks through the fused \texttt{directRecvCopyDirectSend} (\cref{sec:nccl-primitives}), so every step of the FIFO is used.
In slice mode, which \cref{eq:bfly-pos} selects for blocks of $n_cS$ bytes or more, the butterfly moves a window as one strided piece per block (\cref{sec:bfly-res-modes}), and each piece is a call that occupies a whole step even when it is much smaller than a chunk, so most of the steps are mostly empty (\cref{fig:bfly-slots}).
The butterfly also has a barrier between the steps (\cref{sec:bfly-kernels}).
We have not profiled the kernels, so the weight of each factor is not measured, but the first one is the one that the design predicts to matter most.

The dashed lines of \cref{fig:bfly-scaling} show that \texttt{LL} is much worse for \bine{} than for the ring, as for the tree (\cref{sec:tree-ll}); this is consistent with the cost of a call above, since the chunk of \texttt{LL} is much smaller than that of \texttt{SIMPLE} (\cref{fig:nccl-grain}), so a block needs more calls and each of them still costs a full step.
From 256\,kB on and $p\ge16$ \bine{} with \texttt{LL} is $6$--$29\times$ slower than the ring with \texttt{LL}; at 64\,MB and $p\ge8$ \texttt{LL} is $11.5$--$13.4\times$ slower than \texttt{SIMPLE} for \bine{}, against $2.6$--$3.5\times$ for the ring.
As a consequence, \bine{} switches to \texttt{SIMPLE} at much smaller sizes (between 1 and 16\,kB for $p\ge32$, against between 16\,kB and 256\,kB for the ring).

%------------------------------------------------------------------------------
\FloatBarrier
\subsection{The four strategies}\label{sec:bfly-res-modes}

The strategies of \cref{sec:bfly-strategies} differ in where the blocks are placed, but their performance also depends on how the transfers are issued, which is decided by the work partition (\cref{sec:bfly-pos}) and by the cost of a call (\cref{sec:nccl-fifo-cost}).

\paragraph{Slice mode.}
Every channel handles a slice of every block, so a window of $g$ blocks is $g$ strided pieces, and each piece needs its own call (\cref{sec:bfly-fast}).
A call consumes the steps of a full chunk even when it moves much less, so a piece smaller than a chunk wastes most of its step, and a window is sent in $g$ steps instead of the few that a contiguous run of the same size would need (\cref{fig:bfly-slots}).
This is what block-by-block always does, and it is what the other strategies fall back to in slice mode.

\paragraph{Pos mode.}
A channel owns whole blocks, so send and double send move a window as a contiguous run, with full chunks.
The price is that a window is owned by few channels: with $n_c$ channels each one owns $p/n_c$ consecutive positions, and the window of $2^k$ positions of step $k$ intersects only $\lceil 2^k n_c/p\rceil$ of them, the others are idle for that step (\cref{fig:bfly-posuse}).
For small blocks this does not matter, since the time is determined by the latency, and pos mode is selected only there, below the crossover $n_cS$ of \cref{eq:bfly-pos}; for large blocks a window would be moved by a few channels only, and the partition would limit the bandwidth.
Permute cannot use pos mode at all.

The expectation is then that send and double send are much faster than block-by-block where pos mode applies, and that in slice mode they degrade to the calls of block-by-block while keeping their own extra cost: the exchange for send, the permutation for permute.
\Cref{fig:bfly-modes} and \cref{tab:bfly-modes} are in agreement with this.
The partition is the one given by \cref{eq:bfly-pos}: pos mode below $n_cS$ bytes per block, with the default $S=16\,\text{KiB}$, the crossover found experimentally, and slice mode above.
Since $n_c$ was not recorded in the runs, the exact block size of the switch at each point is not known; it lies between 16 and 256\,kB.
% TODO: log the number of channels actually used at each point.

\begin{figure}[!tbp]
\centering
\begin{tikzpicture}[x=1cm, y=1cm, font=\footnotesize,
    slot/.style={draw, minimum height=0.7cm, anchor=south west, inner sep=0, fill=gray!12},
    pc/.style={draw, minimum height=0.7cm, anchor=south west, inner sep=0}]
  \def\w{2.0}% width of a slot = one chunk
  % (a) contiguous run: 4 pieces of 0.4 chunk = 1.6 chunks -> 2 steps
  \node[anchor=west, font=\footnotesize\bfseries] at (0,1.9) {(a) contiguous run (pos mode)};
  \foreach \k in {0,1} {
    \node[slot, minimum width=\w cm] at (2.3*\k,0.9) {};
  }
  \node[pc, fill=ch0!40, minimum width=0.8cm] at (0,0.9) {};
  \node[pc, fill=ch1!40, minimum width=0.8cm] at (0.8,0.9) {};
  \node[pc, fill=ch2!40, minimum width=0.4cm] at (1.6,0.9) {};
  \node[pc, fill=ch2!40, minimum width=0.4cm] at (2.3,0.9) {};
  \node[pc, fill=chunkB!60, minimum width=0.8cm] at (2.7,0.9) {};
  \node[anchor=west] at (5.0,1.25) {2 steps for 4 pieces: only the last one is partly empty};
  % (b) strided pieces: 4 calls, 4 steps
  \node[anchor=west, font=\footnotesize\bfseries] at (0,-0.3) {(b) strided pieces (slice mode)};
  \foreach \k/\c in {0/ch0,1/ch1,2/ch2,3/chunkB} {
    \node[slot, minimum width=\w cm] at (2.3*\k,-1.3) {};
    \node[pc, fill=\c!40, minimum width=0.8cm] at (2.3*\k,-1.3) {};
  }
  \node[anchor=west] at (9.4,-0.95) {4 steps for 4 pieces,};
  \node[anchor=west] at (9.4,-1.45) {each $60\,\%$ empty};
  % legend
  \node[slot, minimum width=0.5cm] at (0,-2.5) {};
  \node[anchor=west] at (0.6,-2.15) {FIFO step (one chunk)};
  \node[pc, fill=ch0!40, minimum width=0.5cm] at (4.2,-2.5) {};
  \node[anchor=west] at (4.8,-2.15) {piece of a block ($0.4$ of a chunk)};
\end{tikzpicture}
\caption{Cost of a call. A call occupies the FIFO steps of a whole chunk, whatever it moves. A window of four blocks whose pieces are $0.4$ of a chunk each is sent in two steps if the pieces are adjacent (a), and in four mostly empty steps if each piece is a separate call (b).}
\label{fig:bfly-slots}
\end{figure}

\begin{figure}[!tbp]
\centering
\begin{tikzpicture}[x=1.05cm, y=0.95cm, font=\footnotesize,
    cell/.style={draw, minimum width=1cm, minimum height=0.8cm, inner sep=0pt}]
  \foreach \row/\n/\act in {0/1/1,1/2/1,2/4/2} {
    \foreach \y/\c in {0/0,1/0,2/1,3/1,4/2,5/2,6/3,7/3} {
      \ifnum\y<\n \edef\fc{pch\c}\def\dc{black}\else\def\fc{white}\def\dc{gray!50}\fi
      \node[cell, fill=\fc, draw=\dc] at (\y,-\row) {};
    }
    \node[anchor=west] at (7.7,-\row) {$\act$ of $4$ channels};
    \node[anchor=east] at (-0.7,-\row) {step $k{=}\row$};
  }
  \foreach \y in {0,...,7} \node[above] at (\y,0.4) {\y};
  \foreach \c in {0,...,3} {
    \fill[pch\c] (2*\c-0.5,1.35) rectangle ++(0.3,0.3);
    \node[anchor=west] at (2*\c-0.15,1.5) {channel \c};
  }
\end{tikzpicture}
\caption{Use of the channels in pos mode, allgather with $p=8$ and four channels, for the rank of position $0$. Each channel owns two positions; the window sent at step $k$ (colored) has $2^k$ positions and involves only the channels that own part of it. In the reduce-scatter the steps are in the opposite order.}
\label{fig:bfly-posuse}
\end{figure}

\begin{figure}[!tbp]
  \bflypair{AllGather_SIMPLE+LL_best_algo_heatmap_BINE_bv-BUTTERFLY.png}{ReduceScatter_SIMPLE+LL_best_algo_heatmap_BINE_bv-BUTTERFLY.png}{0.95}
  \caption{Fastest \bine{} strategy and protocol for each block size and number of ranks (median latency), without the ring. Each cell lists the other combinations with the percentage by which they are slower. B, S, P, D: block-by-block, send, permute and double send; the suffixes S and LL are the protocols.}
  \label{fig:bfly-modes}
\end{figure}

\begin{table}[!tbp]
\centering
\small
\setlength{\tabcolsep}{4pt}
\renewcommand{\arraystretch}{1.1}
\sisetup{table-format=1.2, detect-weight=true, mode=text}
\caption{Latency of send, double send and permute relative to block-by-block, with \texttt{SIMPLE} on $64$ ranks (a value below $1$ means faster than block-by-block).}
\label{tab:bfly-modes}
\begin{tabular}{@{}l *{3}{S} @{\hspace{2em}} *{3}{S}@{}}
\toprule
 & \multicolumn{3}{c}{Allgather} & \multicolumn{3}{c}{Reduce-scatter} \\
\cmidrule(lr){2-4}\cmidrule(l){5-7}
Block size & {send} & {double} & {permute} & {send} & {double} & {permute} \\
\midrule
1\,kB & 0.31 & 0.32 & 0.83 & 0.10 & 0.08 & 0.20 \\
16\,kB & 0.95 & 1.15 & 2.43 & 0.62 & 0.75 & 1.41 \\
256\,kB & 1.09 & 1.93 & 3.38 & 1.04 & 2.34 & 1.79 \\
1\,MB & 1.05 & 2.11 & 4.58 & 0.99 & 2.32 & 1.95 \\
4\,MB & 1.09 & 2.61 & 4.99 & 1.07 & 2.50 & 2.00 \\
64\,MB & 1.03 & 2.41 & 4.95 & 1.06 & 2.33 & 2.02 \\
\bottomrule
\end{tabular}
\end{table}

At 1\,kB block-by-block pays one call per block, and is $3\times$ slower than send and double send in the allgather ($609$ against $191$ and $194$\,\textmu s) and even more in the reduce-scatter ($2227$ against $218$ and $184$\,\textmu s; the block-by-block point is irregular, since on $32$ ranks it takes $213$\,\textmu s).
This is the benefit of pos mode, and it disappears between 16 and 256\,kB, where the switch to slice mode takes place ($n_cS$, \cref{eq:bfly-pos}).
From 256\,kB on send is at most $9\,\%$ slower than block-by-block, i.e., it has the same transfers plus its exchange, whose cost shrinks as $p$ grows (at 64\,MB, $1.2$--$1.3\times$ on $4$ and $8$ ranks and $1.03$--$1.06\times$ on $64$) because it moves one block against the $p-1$ of the butterfly.
Permute adds the in-place permutation of the whole buffer, so on $64$ ranks it is $4.6$--$5\times$ slower than block-by-block in the allgather and $2\times$ in the reduce-scatter from 1\,MB on, and it is the best strategy only at 1\,kB on $16$ ranks or fewer, by at most $12\,\%$.
Double send degrades more than the others: from 256\,kB on it is $1.9$--$2.6\times$ slower than block-by-block, although in slice mode it issues the same transfers, and with \texttt{LL} it is up to $15\times$ slower.
We have not identified the cause, and a profile of this strategy is needed.
% TODO: profile double send (per-step times, per-channel traffic).

In summary, only the pos mode makes the strategies other than block-by-block useful, and only for small blocks: the best \bine{} strategy is send for 1--16\,kB on many ranks, block-by-block from 16\,kB on (within $2\,\%$ of the best from 256\,kB on, with \texttt{SIMPLE}), and permute and double send are never worth their extra cost.
